\begin{Solbook}{1}
 \par \noindent  a) $-1$ \hspace {2cm} b) 1  \hfill \qedEX 
\end{Solbook}
\begin{Solbook}{2}
 \par \noindent  a) (3;0) \hspace {2cm} b) $x\ge -1$  \hfill \qedEX 
\end{Solbook}
\begin{Solbook}{3}
 \par \noindent  $h\approx 52,1$.  \hfill \qedEX 
\end{Solbook}
\begin{Solbook}{4}
 \par \noindent  Tiền điện và tiền nước tháng 11 lần lượt là 809\,400 đồng và 165\,600 đồng.  \hfill \qedEX 
\end{Solbook}
\begin{Solbook}{5}
 \par \noindent  \begin {center} \begin {tikzpicture}[line width=0.6pt] \def \R {2.5} \def \A {5.5} \path (0,0) coordinate (O) (\A ,0) coordinate (M) ($(O)!0.5!(M)$) coordinate (I) ; \path [name path=ciro] (O) circle (\R ); \path [name path=cirom] (I) let \p 1=($(I)-(O)$) in circle ({veclen(\x 1,\y 1)}); \path [name intersections={of=ciro and cirom, by={A,B}}]; \coordinate (C) at ($(A)!2!(O)$); \path [name path=mc] (M)--(C); \path [name intersections={of=ciro and mc,by={C,D}}]; \coordinate (H) at ($(A)!(B)!(C)$); \coordinate (At) at ($(A)!3!(M)$); \coordinate (Ct) at ($(C)!10!(B)$); \coordinate (T) at (intersection of C--Ct and A--At); \draw (O) circle (\R ) (A)--(O)--(B) (C)--(O)--(M)--cycle (A)--(B) (A)--(M) (A)--(I)--(B) (M)--(B) (C)--(T)--(M) (B)--(H) (D)--(B) ; \foreach \x /\g in{O/-180,M/-90,A/60,B/-60,C/-90,D/-90,T/0,I/-90,H/150} \draw [fill=white] (\x ) circle (1.5pt) +(\g :.3) node{$\x $}; \end {tikzpicture} \end {center} \text {} \begin {enumerate} \item Gọi $I$ là trung điểm $OM$. Từ các tam giác vuông $OAM$ và $OBM$, chứng minh điểm $I$ cách đều các điểm $O,A,M,B$, từ đó suy ra 4 điểm $O,A,M,B$ cùng thuộc một đường tròn. \item Chứng minh $\triangle MDB\backsim MBC$, từ đó suy ra tỉ số đồng dạng và có điều phải chứng minh. \item Kéo dài $BC$ cắt $AM$ tại $T$. Chứng minh $M$ là trung điểm $AT$ dựa vào tam giác $ABT$, từ đó sử dụng định lý Thales cho tam giác $CAT$ với $BH\parallel TA$ để suy ra $MC$ đi qua trung điểm $BH$. \end {enumerate}  \hfill \qedEX 
\end{Solbook}
\begin{Solbook}{6}
 \par \noindent  $R=6+4\sqrt {3}$.  \hfill \qedEX 
\end{Solbook}
